% Section: definability

\documentclass[../../../include/open-logic-section]{subfiles}

\begin{document}

\olfileid[fr]{nml}{frd}{def}

\olsection{Définissabilité des classes de cadres}

Bien que les réciproques du \olref[acc]{thm:soundschemas}
échouent pour les modèles, elles valent pour les cadres :
pour chacune des propriétés considérées dans le
\olref[acc]{thm:soundschemas}, si la !!{formula} correspondante
est valide dans un \emph{cadre}, alors la relation d'accessibilité
de ce cadre possède la propriété en question. Les !!{formula}s
considérées \emph{définissent} donc les classes de cadres qui
possèdent ces propriétés.

\begin{defn}
  Si $\mClass{F}$ est une classe de cadres, nous disons que
  $!A$ \emph{définit}~$\mClass{F}$ si et seulement si
  $\mModel{F} \Entails !A$ pour tous les cadres
  $\mModel{F} \in \mClass{F}$, et pour eux seuls.
\end{defn}

Établissons maintenant les résultats complets de
définissabilité pour les cadres.

\begin{thm}\ollabel{thm:fullCorrespondence}
Si la !!{formula} indiquée à droite dans le
\olref[acc]{tab:five} est valide dans un cadre~$\mModel{F}$,
alors $\mModel{F}$ possède la propriété correspondante
indiquée à gauche.
\end{thm}

\begin{proof}
  \begin{enumerate}
  \item Supposons \Ax{D} valide dans
    $\mModel{F} = \tuple{W, R}$, c'est-à-dire
    $\mModel{F} \Entails \Box p \lif \Diamond p$. Soient
    $\mModel{M} = \tuple{W, R, V}$ un modèle fondé
    sur~$\mModel{F}$ et $w \in W$. Il faut trouver un monde
    $v$ tel que~$Rwv$. Supposons qu'il n'en existe aucun.
    Alors, pour toute~$!A$, y compris~$p$,
    $\mSat{M}{\Box !A}[w]$ et
    $\mSat/{M}{\Diamond !A}[w]$. Il s'ensuit
    $\mSat/{M}{\Box p \lif \Diamond p}[w]$,
    contrairement à $\mModel{F} \Entails \Box p \lif \Diamond p$.
  \item Supposons \Ax{T} valide dans $\mModel{F}$,
    c'est-à-dire $\mModel{F} \Entails \Box p \lif p$.
    Soit $w \in W$ quelconque; montrons $Rww$.
    Choisissons $V$ de sorte que, pour tout monde~$u$,
    $u \in V(p)$ si et seulement si $Rwu$
    (pour toute variable $q$ distincte de~$p$, $V(q)$ est
    arbitraire; prenons, par exemple, $V(q) = \emptyset$).
    Posons $\mModel{M} = \tuple{W, R, V}$. Par construction,
    pour tout $u$ tel que $Rwu$, nous avons
    $\mSat{M}{p}[u]$ et donc $\mSat{M}{\Box p}[w]$.
    Or, par hypothèse, $\Box p \lif p$ est vraie en~$w$;
    ainsi $\mSat{M}{p}[w]$. Par définition de~$V$,
    cela n'est possible que si~$Rww$.
  \item Raisonnons par contraposition. Supposons
    $\mModel{F}= \tuple{W, R}$ non symétrique; montrons
    que \Ax{B}, soit $p \lif \Box\Diamond p$, n'y est pas
    valide. Il existe alors $u$, $v \in W$ tels que $Ruv$,
    mais pas $Rvu$. Définissons $V$ par la condition
    suivante : pour tout monde~$w$, $w \in V(p)$ si et
    seulement si $Rvw$ ne vaut pas (et choisissons
    arbitrairement les valeurs des autres
    variables). Soit $\mModel{M} =\tuple{W, R, V}$.
    Par définition de~$V$, $\mSat{M}{p}[w]$ pour tout
    $w$ tel que $Rvw$ ne vaut pas; en particulier,
    $\mSat{M}{p}[u]$, car $Rvu$ ne vaut pas.
    En outre, $Rvw$ si et seulement si $w \notin V(p)$;
    aucun monde~$w$ ne satisfait donc à la fois $Rvw$
    et $\mSat{M}{p}[w]$. Ainsi,
    $\mSat/{M}{\Diamond p}[v]$. Puisque~$Ruv$,
    $\mSat/{M}{\Box\Diamond p}[u]$ également. Par
    conséquent, $\mSat/{M}{p \lif \Box\Diamond p}[u]$,
    et \Ax{B} n'est pas valide dans~$\mModel{F}$.
  \item Supposons \Ax{4} valide dans
    $\mModel{F} = \tuple{W,R}$, autrement dit
    $\mModel{F} \Entails \Box p \lif \Box\Box p$.
    Soient $u$, $v$, $w \in W$ tels que $Ruv$ et $Rvw$;
    montrons $Ruw$. Définissons $V$ de sorte que,
    pour tout monde~$z$, $z \in V(p)$ si et seulement si
    $Ruz$ (et choisissons
    arbitrairement les valeurs des autres variables). Soit
    $\mModel{M} =\tuple{W, R, V}$. Par définition de~$V$,
    $\mSat{M}{p}[z]$ pour tout $z$ tel que $Ruz$,
    et donc $\mSat{M}{\Box p}[u]$. Mais, par hypothèse,
    \Ax{4}, soit $\Box p \lif \Box \Box p$, est vraie
    en~$u$; ainsi $\mSat{M}{\Box \Box p}[u]$.
    Comme $Ruv$ et $Rvw$, il en découle
    $\mSat{M}{p}[w]$. Par définition de~$V$, cela
    n'est possible que si $Ruw$, comme voulu.
  \item Procédons par contraposition. Supposons le cadre
    $\mModel{F} = \tuple{W, R}$ non euclidien et montrons
    qu'il falsifie~\Ax{5}, c'est-à-dire
    $\mModel{F} \Entails/ \Diamond p \lif \Box\Diamond p$.
    Il existe donc des mondes $u$, $v$, $w \in W$ tels que
    $Rwu$ et $Rwv$, mais pas~$Ruv$. Définissons $V$ de
    sorte que, pour tout monde $z$, $z \in V(p)$ si et
    seulement si $Ruz$ ne vaut \emph{pas}. Soit
    $\mModel{M} =\tuple{W, R, V}$. L'hypothèse donne
    $\mSat{M}{p}[v]$ et, puisque $Rwv$,
    $\mSat{M}{\Diamond p}[w]$. En revanche, aucun monde
    $y$ ne satisfait à la fois $Ruy$ et
    $\mSat{M}{p}[y]$; donc
    $\mSat/{M}{\Diamond p}[u]$. Comme $Rwu$, il en
    résulte $\mSat/{M}{\Box\Diamond p}[w]$ : \Ax{5},
    soit $\Diamond p \lif \Box\Diamond p$, échoue
    en~$w$.
  \end{enumerate}
\end{proof}

La preuve pour \Ax{D} diffère des autres cas : elle n'utilise
aucune propriété de l'attribution~$V$. Nous avons en effet montré
que, si $\mSat{M}{\Box p \lif \Diamond p}$, alors la relation
du modèle~$\mModel{M}$ est sérielle. \Ax{D} caractérise donc
également les \emph{modèles} sériels, et pas seulement les cadres
sériels.

\begin{cor}\ollabel{prop:D-serial}
  Tout modèle dans lequel \Ax{D} est vraie est sériel.
\end{cor}

\begin{cor}
Chacune des !!{formula}s indiquées à droite dans le
\olref[acc]{tab:five} définit la classe des cadres qui
possèdent la propriété correspondante indiquée à gauche.
\end{cor}

\begin{proof}
  Dans le \olref[acc]{thm:soundschemas}, nous avons prouvé que,
  si la relation d'un modèle possède la propriété indiquée
  à gauche, la !!{formula} correspondante est vraie dans ce
  modèle. Ainsi, si un cadre~$\mModel{F}$ possède cette
  propriété, la !!{formula} correspondante est valide
  dans~$\mModel{F}$. Le
  \olref{thm:fullCorrespondence} établit la réciproque :
  si la !!{formula} indiquée à droite est valide
  dans~$\mModel{F}$, ce cadre possède la propriété
  indiquée à gauche.
\end{proof}

\begin{prob}
Montrez que, si une !!{formula} indiquée à droite dans le
\olref[nml][frd][acc]{tab:anotherfive} est valide dans un
cadre~$\mModel{F}$, alors ce cadre possède la propriété
correspondante indiquée à gauche. Pour cela, considérez un
cadre qui ne possède \emph{pas} cette propriété et
choisissez une attribution~$V$ telle que la !!{formula}
correspondante soit fausse en un monde.
\end{prob}

Le \olref{thm:fullCorrespondence} montre aussi que les propriétés
peuvent se combiner : si \Ax{B} et \Ax{4} sont toutes deux
valides dans~$\mModel{F}$, par exemple, ce cadre est à la fois
symétrique et transitif. De nombreuses logiques modales
importantes sont constituées des !!{formula}s valides dans tous
les cadres réunissant certaines propriétés. On peut donc
également les caractériser comme les ensembles de !!{formula}s
valides dans tous les cadres où les formules qui définissent
ces propriétés sont valides. Ainsi, le système classique
\Log{S4} est l'ensemble de toutes les !!{formula}s valides dans tous
les cadres réflexifs et transitifs, c'est-à-dire dans tous
les cadres où \Ax{T} et~\Ax{4} sont valides. Le système
\Log{S5} est l'ensemble de toutes les formules valides dans tous les
cadres réflexifs, symétriques et euclidiens, c'est-à-dire
dans ceux où \Ax{T}, \Ax{B} et \Ax{5} sont valides.

Les relations logiques entre propriétés de $R$ se reflètent
généralement dans les relations entre les !!{formula}s qui
les définissent. Par exemple, toute relation réflexive est
sérielle : si \Ax{T} est valide dans un cadre, \Ax{D} y est
donc valide aussi. (Il ne s'agit \emph{pas} ici d'une relation
de conséquence logique entre ces formules : il peut arriver que
$\mSat{M}{\Box p \lif p}[w]$ sans que
$\mSat{M}{\Box p \lif \Diamond p}[w]$.) Voici quelques
relations de ce type.

\begin{prop}\ollabel{prop:relation-facts}
  Soit $R$ une relation binaire sur un ensemble $W$. Alors :
  \begin{enumerate}
  \item Si $R$ est réflexive, elle est sérielle.
  \item Si $R$ est symétrique, elle est transitive si et
    seulement si elle est euclidienne.
  \item Si $R$ est symétrique ou euclidienne, elle est
    faiblement dirigée (elle possède la « propriété du
    diamant »).
  \item Si $R$ est euclidienne, elle est faiblement connexe.
  \item Si $R$ est fonctionnelle, elle est sérielle.
  \end{enumerate}
\end{prop}

\begin{prob}
  Démontrez la \olref[nml][frd][def]{prop:relation-facts}.
\end{prob}
\begin{editorial}Dans le cas de la sérialité, la source omettait un indice de monde dans l’assertion portant sur la nécessité ; il est rétabli, car l’argument porte sur le monde choisi. Une parenthèse du cas de la réflexivité est aussi replacée hors de l’expression mathématique. Les occurrences de satisfaction portant sur les noms des axiomes sont développées avec leurs formules, sans changer les conclusions.\end{editorial}
\end{document}
